% Complete Section 7.1, original PDF 92-93 / printed 82-83.
% block: complete-polynomial-forms-heading page:92 type:section
\subsection{\BT{The complete polynomial spaces of differential forms}{微分形式的完全多项式空间}}\label{sec:complete-polynomial-forms}

% block: finite-element-assembly-definition page:92 type:ordinary-paragraph
\BT{The spaces $\mathcal P_r\Lambda^k(\mathcal T_h)$ and $\mathcal P_r^-\Lambda^k(\mathcal T_h)$ will be constructed through the usual process of \emph{finite element assembly}. In general, a finite element space is determined by the specification of the triangulation $\mathcal T_h$, and for each element $T\in\mathcal T_h$:}{空间 $\mathcal P_r\Lambda^k(\mathcal T_h)$ 和 $\mathcal P_r^-\Lambda^k(\mathcal T_h)$ 将通过通常的\emph{有限元组装}过程构造. 一般来说, 指定剖分 $\mathcal T_h$, 并对每个单元 $T\in\mathcal T_h$ 指定下列内容, 就确定了一个有限元空间:}
\begin{itemize}
\item\BT{a finite dimensional space of shape functions $V(T)$ on $T$, and}{$T$ 上的有限维形函数空间 $V(T)$;}
\item\BT{a set of degrees of freedom (DOFs) for $V(T)$ (i.e., a basis for its dual space), with each DOF associated to a particular face of $T$.}{$V(T)$ 的一组自由度 (DOFs), 即其对偶空间的一组基, 其中每个自由度与 $T$ 的一个特定面关联.}
\end{itemize}

% block: finite-element-assembled-space page:92 type:ordinary-paragraph
\BT{The assembled finite element space is then defined to consist of all functions on $\Omega$ whose restriction to each $T$ belongs to the shape function space $V(T)$ and for which the DOFs are single-valued (in the sense that when two elements share a common face, the corresponding DOFs on the face take the same value).}{组装后的有限元空间定义为 $\Omega$ 上所有满足下列条件的函数: 在每个 $T$ 上的限制属于形函数空间 $V(T)$, 且自由度是单值的 (即两个单元共享一个面时, 面上的相应自由度取相同值).}

% block: complete-polynomial-shape-functions pages:92,93 type:ordinary-paragraph
\BT{In the present case, the shape function spaces will be spaces of differential $k$-forms with polynomial coefficients. For $\mathcal P_r\Lambda^k(\mathcal T_h)$, the choice of shape function space $V(T)$ is an obvious one, namely, the complete polynomial space}{在当前情形中, 形函数空间是多项式系数微分 $k$-形式的空间. 对 $\mathcal P_r\Lambda^k(\mathcal T_h)$, 形函数空间 $V(T)$ 的选取是显然的, 即完全多项式空间}
\[
\mathcal P_r\Lambda^k(T)=\left\{\,\sum_{1\leq\sigma_1<\cdots<\sigma_k\leq n}p_\sigma\,dx^\sigma\ \middle|\ p_\sigma\in\mathcal P_r(T)\,\right\}
\]
\BT{of a differential $k$-form with polynomial coefficients of degree at most $r$. Its dimension is easily calculated:}{它由系数为次数不超过 $r$ 的多项式的微分 $k$-形式构成. 其维数容易计算:}
\begin{equation}\label{eq:complete-polynomial-form-dimension}
\begin{aligned}
\dim\mathcal P_r\Lambda^k(T)&=\binom nk\times\dim\mathcal P_r(T)\\
&=\binom nk\binom{n+r}n=\binom{r+n}{r+k}\binom{r+k}r.
\end{aligned}
\end{equation}

% block: polynomial-de-rham-complex page:93 type:ordinary-paragraph
\BT{Note that $d\mathcal P_r\Lambda^k\subset\mathcal P_{r-1}\Lambda^{k+1}$, i.e., the exterior derivative lowers the polynomial degree at the same time as it raises the form degree. Therefore, for each $r$ we have a polynomial subcomplex of the de Rham complex:}{注意 $d\mathcal P_r\Lambda^k\subset\mathcal P_{r-1}\Lambda^{k+1}$, 即外微分在提高形式次数的同时降低多项式次数. 因此, 对每个 $r$, 都有 de Rham 复形的一个多项式子复形:}
\[
0\to\mathcal P_r\Lambda^0\xrightarrow{d}\mathcal P_{r-1}\Lambda^1\xrightarrow{d}\cdots\xrightarrow{d}\mathcal P_{r-n}\Lambda^n\to0.
\]
\BT{(We understand $\mathcal P_s\Lambda^k$ to be $0$ if $s<0$.) Moreover, as proved below in Corollary~\ref{cor:polynomial-de-rham-exactness}, this complex is exact, except for the constant functions in $\mathcal P_r\Lambda^0$, which form the kernel of the first $d$. That is, if $\omega\in\mathcal P_s\Lambda^k$ for some $k>0$ and $d\omega=0$, then $\omega=d\mu$ for some $\mu\in\mathcal P_{s+1}\Lambda^{k-1}$. The proof of this fact will involve a very valuable tool called the \emph{Koszul complex}. The same tool will be used to define DOFs and to define an alternative space of shape functions.}{(若 $s<0$, 约定 $\mathcal P_s\Lambda^k$ 为 $0$.) 此外, 如下面 Corollary~\ref{cor:polynomial-de-rham-exactness} 所证明, 除 $\mathcal P_r\Lambda^0$ 中构成第一个 $d$ 的核的常函数外, 此复形是正合的. 即若对某个 $k>0$ 有 $\omega\in\mathcal P_s\Lambda^k$, 且 $d\omega=0$, 则对某个 $\mu\in\mathcal P_{s+1}\Lambda^{k-1}$ 有 $\omega=d\mu$. 证明这一事实将涉及称为 \emph{Koszul 复形}的非常有用的工具. 同一工具还将用于定义自由度, 以及定义另一种形函数空间.}
